\documentclass[letterpaper,11pt]{article}

\usepackage{latexsym}
\usepackage[empty]{fullpage}
\usepackage{titlesec}
\usepackage{marvosym}
\usepackage[usenames,dvipsnames]{color}
\usepackage{verbatim}
\usepackage{enumitem}
\usepackage[hidelinks]{hyperref}
\usepackage{fancyhdr}
\usepackage[english]{babel}
\usepackage{tabularx}
\usepackage{fontawesome5}
\usepackage{multicol}
\setlength{\multicolsep}{-3.0pt}
\setlength{\columnsep}{-1pt}
\input{glyphtounicode}


\pagestyle{fancy}
\fancyhf{} % clear all header and footer fields
\fancyfoot{}
\renewcommand{\headrulewidth}{0pt}
\renewcommand{\footrulewidth}{0pt}

% Adjust margins
\addtolength{\oddsidemargin}{-0.5in}
\addtolength{\evensidemargin}{-0.5in}
\addtolength{\textwidth}{1in}
\addtolength{\topmargin}{-.7in}
\addtolength{\textheight}{1.2in}

\urlstyle{same}

\raggedbottom
\raggedright
\setlength{\tabcolsep}{0in}

% Sections formatting
\titleformat{\section}{
  \vspace{-4pt}\scshape\raggedright\large\bfseries
}{}{0em}{}[\color{black}\titlerule \vspace{-5pt}]

% Ensure that generated pdf is machine readable/ATS parsable
\pdfgentounicode=1

%-------------------------
% Custom commands
\newcommand{\resumeItem}[1]{
  \item\small{
    {#1 \vspace{-4pt}}
  }
}

\newcommand{\classesList}[4]{
    \item\small{
        {#1 #2 #3 #4 \vspace{-4pt}}
  }
}

\newcommand{\resumeSubheading}[4]{
  \vspace{-2pt}\item
    \begin{tabular*}{1.0\textwidth}[t]{l@{\extracolsep{\fill}}r}
      \textbf{#1} & \textbf{\small #2} \\
      \textit{\small#3} & \textit{\small #4} \\
    \end{tabular*}\vspace{-8pt}
}

\newcommand{\eduSubheading}[2]{
  \vspace{-2pt}\item
    \begin{tabular*}{1.0\textwidth}[t]{l@{\extracolsep{\fill}}r}
      \textbf{#1} & \textbf{\small #2} \\
    \end{tabular*}\vspace{-8pt}
}

\newcommand{\resumeSubSubheading}[2]{
    \item
    \begin{tabular*}{0.9\textwidth}{l@{\extracolsep{\fill}}r}
      \textit{\small#1} & \textit{\small #2} \\
    \end{tabular*}\vspace{-8pt}
}

\newcommand{\resumeProjectHeading}[2]{
    \item
    \begin{tabular*}{1\textwidth}{l@{\extracolsep{\fill}}r}
      \small#1 & \textbf{\small #2}\\
    \end{tabular*}\vspace{-8pt}
}

\newcommand{\resumeSubItem}[1]{\resumeItem{#1}\vspace{-4pt}}

\renewcommand\labelitemi{$\vcenter{\hbox{\tiny$\bullet$}}$}
\renewcommand\labelitemii{$\vcenter{\hbox{\tiny$\bullet$}}$}

\newcommand{\resumeSubHeadingListStart}{\begin{itemize}[leftmargin=0.0in, label={}]}
\newcommand{\resumeSubHeadingListEnd}{\end{itemize}}
\newcommand{\resumeItemListStart}{\begin{itemize}}
\newcommand{\resumeItemListEnd}{\end{itemize}\vspace{-5pt}}

%-------------------------------------------
%%%%%%  RESUME STARTS HERE  %%%%%%%%%%%%%%%%%%%%%%%%%%%%


\begin{document}

\begin{center}
    {\Huge \scshape Akash Jaiswal} \\ \vspace{4pt}

    \small Platform Engineer II @ Oracle | Platform Engineering, Cloud Native \& AI Systems \\ \vspace{3pt}

    \href{mailto:akashjaiswal3846@gmail.com}
    {\raisebox{-0.2\height}\faEnvelope\ \underline{akashjaiswal3846@gmail.com}} ~
    \href{https://linkedin.com/in/akashjaiswal03}
    {\raisebox{-0.2\height}\faLinkedin\ \underline{/akashjaiswal03}} ~
    \href{https://github.com/jaiakash}
    {\raisebox{-0.2\height}\faGithub\ \underline{/jaiakash}} ~
    \href{https://twitter.com/a_kashhhhhh_}
    {\raisebox{-0.2\height}\faTwitter\ \underline{/akashhhhhh}}

    \vspace{-8pt}
\end{center}


%-----------EXPERIENCE-----------
\section{Experience}
  \resumeSubHeadingListStart
    \resumeSubheading
    {Oracle India Pvt Ltd}{June 2024 -- Present}
    {Platform Engineer II}{Hyderabad, Telangana}
    \resumeItemListStart

    \resumeItem{Part of the \textbf{Fusion Release Engineering} organization, building and operating Kubernetes-native CI/CD platforms and developer infrastructure serving \textbf{5,000+ developers} across \textbf{1,200+ projects}; developed platform services using \textbf{Kubernetes, OCI, Podman, Java (Helidon), Node.js, Oracle JET, and Oracle Database}.}
    
    \resumeItem{Built core components of \textbf{Ocean}, a Kubernetes-native CI platform replacing legacy Jenkins infrastructure; implemented workload scheduling, autoscaling, and resource quota management to optimize cluster utilization and support thousands of concurrent CI workloads across Oracle Fusion Applications.}
    
    \resumeItem{Developed \textbf{Darwin ITS}, a self-service developer platform enabling pipeline simulation, automated validation, and AI-assisted workflow analysis prior to production deployment; adopted by \textbf{50+ developers} and significantly reducing validation turnaround time.}
    
    \resumeItemListEnd

\resumeSubheading
    {\href{https://summerofcode.withgoogle.com/programs/2025/projects/fwZkvPr0}{Kubeflow, CNCF (Google Summer of Code)}}{April 2025 -- Present}
    {Open Source Contributor \& Mentor | GSoC 2025, 2026}{Remote}
    \resumeItemListStart
    
    \resumeItem{\textbf{GSoC 2025 Contributor}: Designed and deployed GPU-enabled Kubernetes infrastructure for Kubeflow's LLM Blueprint CI on Oracle Kubernetes Engine (OKE); provisioned clusters, configured GitHub self-hosted runners, and implemented approval-gated CI workflows to enable reliable testing of AI/LLM workloads before upstream release.}

    \resumeItem{Serving as \textbf{OCI Tenancy Administrator} for Kubeflow community infrastructure, owning deployment and operations of GPU-enabled Kubernetes clusters on OCI; managing IAM, networking, Argo CD-based GitOps workflows, secrets, contributor access, and platform scalability for AI/ML workloads.}
    
    \resumeItem{\textbf{GSoC 2026 Mentor}: Co-mentoring projects focused on \textbf{MCP integration} and \textbf{OpenTelemetry observability} for Kubeflow SDK, enabling LLM agents to manage TrainJob lifecycles and improving distributed tracing, metrics, and operational visibility across AI/ML workflows.}

\resumeItemListEnd


\resumeSubheading
{Selected Projects}{2022 -- Present}
{AI Infrastructure \& Open Source}{}

\resumeItemListStart

\resumeItem{\textbf{Self-Hosted LLM Platform}: Built and operated a GPU-backed LLM serving platform on OCI using \textbf{vLLM}, \textbf{Open WebUI}, NGINX, and OpenAI-compatible APIs, enabling secure self-hosted inference workloads on GPU infra.}

\resumeItem{\textbf{\href{https://www.linkedin.com/feed/update/urn:li:activity:6933613937011617792/}{Google Summer of Code 2022 (CCExtractor)}}: Built and deployed a cross-platform Flutter application and Node.js backend that allowed users to flag urgent calls and deliver priority notifications, enabling interruption-aware communication workflows on mobile devices.}


\resumeItemListEnd

\vspace{-10pt}


\section{Open Source Leadership \& Community}
    \resumeItemListStart
    
    \resumeItem{
    \textbf{Open Source Contributions}: Contributed to \textbf{Kubeflow}, OWASP, Keploy, Oppia, Processing Foundation, and other open-source projects across cloud-native, AI/ML, security, and developer tooling ecosystems.
    }
    
    \resumeItem{
    \textbf{Conference Speaker}: 
    \href{https://ossindia2026.sched.com/event/2KNF7}{\textbf{Open Source Summit India 2026}},
    \href{https://gdg.community.dev/events/details/google-gdg-cloud-nagpur-presents-techfusion-summit-2026/}{\textbf{TechFusion Summit 2026 (GDG Cloud Nagpur)}},
    \textbf{CNCF Hyderabad Meetup}.
    }
    
    \resumeItem{
    \textbf{Community Leadership}: Served as \textbf{Technical Secretary} at NIT Trichy, overseeing 17 technical clubs and Pragyan (budget: INR 50+ lakhs); founded and scaled the \textbf{TeCOS} open-source community to \textbf{1,000+ members}, mentoring students in open source and GSoC.
    }
    
    \resumeItem{
    \textbf{Community Recognition}: Awarded scholarships by the \textbf{Linux Foundation}, \textbf{CNCF}, and \textbf{Wikimedia Foundation} to attend international open-source conferences including
    \href{https://www.linkedin.com/posts/akashjaiswal03_opensourcesummit-openssf-opensearchcon-ugcPost-7392580382480658433-a8-J}{\textbf{Open Source Summit Korea 2025}},
    \href{https://blog.kubeflow.org/kubecon/community/2025/08/23/kubecon-2025-india-kubeflow.html}{\textbf{KubeCon India 2025}},
    \href{https://www.cncf.io/blog/2025/03/17/kubecon-cloudnativecon-india-2024-from-dreams-to-reality-journey/}{\textbf{KubeCon India 2024}},
    and
    \href{https://indicwiki.iiit.ac.in/summit-2023/}{\textbf{Wikimedia Summit 2023}}.
    }
    
    \resumeItemListEnd


%-----------PROGRAMMING SKILLS-----------
\section{Technical Skills}
\begin{itemize}[leftmargin=0.15in, label={}]
\small{\item{

\textbf{Languages}{: Go, Python, Java, JavaScript, TypeScript} \\
\textbf{Platform Engineering \& DevOps}{: Kubernetes, Docker/Podman, Helm, Argo CD, GH Actions, Terraform, CI/CD} \\
\textbf{AI/ML Infrastructure}{: Kubeflow, MCP, GPU Infrastructure, LLM Deployment, vLLM} \\
\textbf{Cloud Platforms}{: Oracle Cloud Infrastructure (OCI), AWS,  GCP} \\
\textbf{Observability}{: OpenTelemetry, Prometheus, Grafana} \\
\textbf{Databases}{: Oracle Database, PostgreSQL, MongoDB} \\
\textbf{Certifications}{: CKAD, AI/ML Toolkits with Kubeflow (LFS147)}

}}
\end{itemize}
 \vspace{-20pt}

%-----------EDUCATION-----------
\section{Education}
  \resumeSubHeadingListStart
    \resumeSubheading
      {National Institute of Technology, Tiruchirappalli}{Nov 2020 -- June 2024}
      {Bachelor of Technology in Civil Engineering (Minor: Computer Science) | CGPA: 8.1/10}{Tamil Nadu, India}
  \resumeSubHeadingListEnd
  
\end{document}